\section{Выводы к Главе 1}\label{sec:review_conclusion}

\begin{enumerate}
	\item При распространении мощных синхронизированных оптических импульсов по маломодовым ОВ возможно развитие н-о процессов третьего порядка, таких, как ВКР, ЧВС. Вопрос о взаимном влиянии ЧВС различных типов в литературе исследован слабо и является открытым.
	
	\item Эффективная ГВГ в германосиликатных ОВ обусловлена формированием $\chi^{(2)}$-решётки в процессе оптического полинга. КФС измерялись лишь в стационарном состоянии, а их временная эволюция в процессе записи решётки не исследовалась. Изучение этой эволюции позволит сопоставить предсказания моделей с экспериментом и выявить их ограничения.
	
	\item При измерениях коэффициента оптического поглощения н-о кристаллов метод пьезоэлектрической резонансной лазерной калориметрии (ПРЛК) позволяет измерять усреднённую температуру кристалла, исключая погрешности, связанные с внешним термодатчиком. Однако стандартные подходы к обработке кинетик предполагают постоянство температуры окружающей среды и коэффициента теплоотдачи. Отсутствуют поправочные выражения, снижающие влияние этих факторов.
\end{enumerate}

